\subsection{Evaluation Settings}\label{sec:setup}
\emph{Hardware and software.} Apple M4 (16\,GB), macOS~26.3, Python~3.14 with NumPy~2.4, scikit-learn~1.8 and faiss-cpu~1.13.2, C code compiled with clang~17 (\texttt{-O3}), PostgreSQL~16.14 with \texttt{pgvector}~0.8.7. Every latency is measured on one thread; each plan runs once as a warm-up and twice timed, and we report the slower run. PostgreSQL uses \texttt{shared\_buffers}=4\,GB, no parallel query or JIT, and \texttt{pgvector}'s default HNSW parameters; its latencies are client wall-clock times, including planning.

\emph{Datasets.} \textbf{ARXIV-384}: one million papers from the CC0 arXiv-abstracts-2021 dump~\cite{arxivdump}, embedded with all-MiniLM-L6-v2~\cite{reimers2019sbert} (384 dimensions, title and abstract), with real attributes (archive, published, year, numbers of categories, authors and abstract words) and a foreign key to its 172 primary categories. \textbf{AMAZON-384}: $990{,}000$ products from Amazon Reviews 2023~\cite{hou2024amazon} in four departments, embedded the same way, with real attributes (department, categories, price, rating, number of ratings, features, description length) and a foreign key to about $144{,}000$ stores. Its filters are built from the query product, as a shopper would filter (same or another department, a price band, a minimum rating, a seller tier). \textbf{SYN-128} and \textbf{SYN-768}: one million synthetic vectors from a topic mixture, with attributes generated from the topics at graded strength. Query vectors are held-out items (real data) or noisy data vectors (synthetic). Both real benchmarks are public with exact answers~\cite{hvobench}.

\emph{Workloads.} Filters combine 1--3 conditions: ranges, \texttt{IN} lists and conditions on the linked table. On SYN-128, SYN-768 and ARXIV-384 they are drawn independently of the query, so filter and query are related only through the data itself. Each benchmark has $1000$ training and $1500$ test queries (queries with no match are dropped); the test splits contain $\nAntiC/\nIndepC/\nPosC$ (ARXIV-384) and $\nAntiD/\nIndepD/\nPosD$ (AMAZON-384) anti-correlated/independent/positively correlated queries. We use $k=10$ unless stated. Everything fitted uses only the training split.

\emph{Baselines.} Estimators: PostgreSQL-style 1-D histograms with independence; the exact global count (Milvus-style); a $1\%$ sample; and an Exqutor-style HNSW probe of the 256 nearest neighbors. Each feeds the same cost model. Routers also include the three single plans and the best global-selectivity rules (two- and three-way thresholds tuned on the training split). Inside PostgreSQL we compare \texttt{pgvector} with default settings, \texttt{pgvector}'s iterative scan (\texttt{strict\_order}), Exqutor~\cite{exqutor2026} built from its released code, both forced plans, and the HVO hook with and without the recall constraint. Exqutor's estimator only activates for vector \emph{range} predicates over more than two tables, so on these queries it plans like \texttt{pgvector}~0.7.1.

\emph{Metrics and statistics.} The oracle runs every plan and takes the fastest with Recall@$k\ge0.9$; other routers are charged their plan's latency plus their planning time. Regret is a router's latency divided by the oracle's. For differences between routers we report paired 95\% percentile-bootstrap intervals over test queries (2000 resamples); they capture variation across queries in one run per setting, not repeated-run timing noise.

